\chapter{Desglose de \texttt{src/universidad-proyecto}}

\section{Resumen funcional}

La carpeta \texttt{src/universidad-proyecto} contiene la solución aplicada completa: cargadores de datos, construcción del grafo de conflictos, solucionadores, API y frontend. El flujo principal de ejecución es:

\begin{enumerate}[leftmargin=1.5cm]
  \item Carga y normalización de datos (\texttt{scheduler-data-loader.ts}).
  \item Construcción del grafo de conflictos (\texttt{ConflictGraphBuilder}).
  \item Separación por tipo de salón y ejecución de algoritmos de coloreo (\texttt{ScheduleSolver}).
  \item Generación del horario semanal final (mapeo vértice→salón→franja).
  \item Exposición por API (controladores y rutas) y visualización en frontend.
\end{enumerate}

\section{Punto de entrada}

El archivo \texttt{index.ts} solo arranca la API:

\begin{lstlisting}[language=TypeScript,caption={Entrada principal del módulo universitario}]
import {SchedulerApiServer} from "./api/server";

const port = parseInt(process.env.PORT || "3002", 10);
const server = new SchedulerApiServer(port);
server.start();
\end{lstlisting}

\section{Componentes clave}

\subsection{Cargador de datos}
Normaliza JSON de entrada en estructuras tipadas (materias, salones, franjas) para el resto del pipeline.

\subsection{ConflictGraphBuilder}
Convierte materias y grupos en vértices y añade aristas cuando existe conflicto (misma franja, docente, semestre compartido, tipo incompatible).

\subsection{ScheduleSolver}
Recibe el grafo de conflictos y el listado de salones. Separa subgrafos por tipo de salón y aplica el algoritmo de coloreo seleccionado. Devuelve un mapeo de vértices a salones.

\begin{lstlisting}[language=TypeScript,caption={Ejecución del solucionador global}]
const input = SchedulerDataLoader.build().getData();
const { graph } = new ConflictGraphBuilder(input).build();
const algoritmo = new DSatur<GrupoData>(graph, input.salones.map(s => s.id));
const horario = new ScheduleBuilder(graph, input.salones).buildHorario(algoritmo);
\end{lstlisting}

\subsection{ScheduleService}
Singleton que mantiene en memoria el resultado del cálculo global para evitar recomputar en cada solicitud. Provee métodos para obtener \texttt{getHorarioGlobal()} y \texttt{getGrafoInfo()}.

\begin{lstlisting}[language=TypeScript,caption={Servicio global en memoria}]
async getHorarioGlobal(): Promise<HorarioSemanal> {
  if (this.horarioGlobal) return this.horarioGlobal;

  const input = SchedulerDataLoader.build().getData();
  const { graph } = new ConflictGraphBuilder(input).build();
  const algoritmo = new DSatur<GrupoData>(graph, input.salones.map(s => s.id));
  const horario = new ScheduleBuilder(graph, input.salones).buildHorario(algoritmo);

  this.horarioGlobal = horario;
  return horario;
}
\end{lstlisting}

\subsection{StudentScheduleSolver}
Genera horarios personalizados para un estudiante a partir del horario global y los grupos elegidos. Integra Dijkstra y backtracking según disponibilidad.

\begin{lstlisting}[language=TypeScript,caption={Creación del horario estudiantil}]
async createSchedule(horario: HorarioSemanal, salones: SalonInput[], request: StudentScheduleRequest) {
  const solver = new StudentScheduleSolver(request);
  const resultado = solver.execute();
  this.horariosEstudiantes.set(request.estudianteId, resultado);
  return resultado;
}
\end{lstlisting}

\section{Rutas y controladores principales}

Rutas expuestas:
\begin{itemize}[leftmargin=1.5cm]
  \item GET /schedule
  \item GET /schedule/grafo
  \item POST /student/schedule
  \item GET /student/schedule/:estudianteId
  \item GET /student/schedules
\end{itemize}

\section{Consideraciones de diseño}

\begin{itemize}[leftmargin=1.5cm]
  \item Separación clara entre construcción del grafo (lógica) y la API (exposición).
  \item Caché en memoria para mejorar latencia en consultas repetidas.
  \item Modularidad para poder reemplazar algoritmos de coloreo sin tocar la capa de datos.
\end{itemize}
